% !TEX TS-program = xelatex
% !TEX encoding = UTF-8 Unicode
% -*- coding: UTF-8; -*-
% vim: set fenc=utf-8

%%%%%%%%%%%%%%%%%%%%%%%%%%%%%%%%%%%%%%%%%%%%%%%%%%%%%%%%%%%%%%%%%
%% CV.tex
%% <https://github.com/zachscrivena/simple-resume-cv>
%% This is free and unencumbered software released into the
%% public domain; see <http://unlicense.org> for details.
%%%%%%%%%%%%%%%%%%%%%%%%%%%%%%%%%%%%%%%%%%%%%%%%%%%%%%%%%%%%%%%%%

% See "README.md" for instructions on compiling this document.

\documentclass[letterpaper,MMMyyyy,nonstopmode]{simpleresumecv}
\usepackage{fontawesome5}
\usepackage{xcolor}
% Class options:
% a4paper, letterpaper, nonstopmode, draftmode
% MMMyyyy, ddMMMyyyy, MMMMyyyy, ddMMMMyyyy, yyyyMMdd, yyyyMM, yyyy

%%%%%%%%%%%%%%%%%%%%%%%%%%%%%%%%%%%%%%%%%%%%%%%%%%%%%%%%%%%%%%%%%
%% PREAMBLE.
%%%%%%%%%%%%%%%%%%%%%%%%%%%%%%%%%%%%%%%%%%%%%%%%%%%%%%%%%%%%%%%%%

% CV Info (to be customized).
\newcommand{\CVAuthor}{Zhenrui Zheng}
\newcommand{\CVTitle}{Zhenrui Zheng's CV}
\newcommand{\CVNote}{CV compiled on {\today}}
\newcommand{\CVWebpage}{https://scholar.google.com/citations?user=KPpd1pYAAAAJ}


% PDF settings and properties.
\hypersetup{
pdftitle={\CVTitle},
pdfauthor={\CVAuthor},
pdfsubject={\CVWebpage},
pdfcreator={XeLaTeX},
pdfproducer={},
pdfkeywords={},
unicode=true,
bookmarks=true,
bookmarksopen=true,
pdfstartview=FitH,
pdfpagelayout=OneColumn,
pdfpagemode=UseOutlines,
hidelinks,
breaklinks}

% Shorthand.
\newcommand{\Code}[1]{\mbox{\textbf{#1}}}
\newcommand{\CodeCommand}[1]{\mbox{\textbf{\textbackslash{#1}}}}

%%%%%%%%%%%%%%%%%%%%%%%%%%%%%%%%%%%%%%%%%%%%%%%%%%%%%%%%%%%%%%%%%
%% ACTUAL DOCUMENT.
%%%%%%%%%%%%%%%%%%%%%%%%%%%%%%%%%%%%%%%%%%%%%%%%%%%%%%%%%%%%%%%%%

\begin{document}

%%%%%%%%%%%%%%%
% TITLE BLOCK %
%%%%%%%%%%%%%%%

\Title{\CVAuthor}

\begin{SubTitle}
\faIcon{envelope} \texttt{ch4acko3@outlook.com}
\,\SubBulletSymbol\,
\faIcon{mobile} \texttt{+86\,159-7977-6837}
\,\SubBulletSymbol\,
\href{\CVWebpage}
{\faIcon{graduation-cap} Google Scholar Profile}
\end{SubTitle}

\begin{Body}

%%%%%%%%%%%%%%%
%% EDUCATION %%
%%%%%%%%%%%%%%%

\Section
{Education}
{Education}
{PDF:Education}

\Entry
\textbf{The Chinese University of Hong Kong, Shenzhen}

\Gap
\BulletItem
Ph.D. Candidate in Computer Science | Expected Graduation: June 2030
\hfill
\DatestampYMD{2025}{9}{1} --
Present
\BulletItem
Advisor: Prof. Chenjun Xiao
\SubBulletItem
Research Interests: Applying Machine Learning to Sequential Decision-Making Tasks (e.g., Reasoning), Generative Models and Reinforcement Learning. I am particularly driven by the pursuit of elegant, robust, and theoretically sound solutions to fundamental problems in research.

\BigGap

\Entry
\textbf{Harbin Engineering University}

\Gap
\BulletItem
B.Eng. in Aircraft Design and Engineering
\hfill
\DatestampYMD{2021}{9}{15} --
\DatestampYMD{2025}{6}{15}
\begin{Detail}
\SubBulletItem
Weighted average score: 86.93/100 (Ranked 18/65)
% \SubBulletItem
% Mathematics course grades: Calculus(94), Linear Algebra(95), Probability and Statistics(95)
% \SubBulletItem
% English test scores: CET-4(641), CET-6(565), IELTS(7.0)
\end{Detail}



%%%%%%%%%%%%%%%%%%%%%%%%%%%
%% AWARDS & SCHOLARSHIPS %%
%%%%%%%%%%%%%%%%%%%%%%%%%%%

\Section
{Awards \&\newline
Scholarships}
{Awards \& Scholarships}
{PDF:AwardsAndScholarships}

\BulletItem
\href{https://neurips.cc/virtual/2024/competition/84793}{\textbf{NeurIPS Competition 2024: AIGB Track Winner}} (Final standing 4/793)
\hfill
\DatestampYM{2024}{11}
\begin{Detail}
\Item
Auto-Bidding in Large-Scale Auctions: Learning Decision-Making in Uncertain and Competitive Games
\Item
Implement offline reinforcement-learning agents based on generative models to solve sequential decision-making and game problems in advertising. Compete with professional algorithm engineers from well-known companies/universities and get fourth place (Other listed teams: Kuaishou\textsuperscript{1st}, BUAA\textsuperscript{2nd}, ByteDance\textsuperscript{3rd}, \textbf{My Team}\textsuperscript{4th},  PKU\textsuperscript{10th}, THU\textsuperscript{15th}, Chinese Academy of Sciences\textsuperscript{16th}, SNU\textsuperscript{22th})
\end{Detail}

\Gap
\BulletItem
\href{https://icpc.global/regionals/finder/ICPC-Hangzhou-2023}{\textbf{The 2022 ICPC Asia Hangzhou Regional Contest}}: Silver Medal
\hfill
\DatestampYMD{2022}{12}{04}
\begin{Detail}
\Item
The International Collegiate Programming Contest (ICPC) is an annual multi-tiered competitive programming competition among the universities of the world. Directed by ICPC Executive Director and Baylor Professor William B. Poucher, the ICPC operates autonomous regional contests covering six continents.
\end{Detail}
% \begin{Detail}
% \Item
% For outstanding scientific contributions in the fields of lasers and climate change.
% \end{Detail}

% \Gap
% \BulletItem
% The 2022 ICPC Asia Hefei Regional Contest: Bronze Medal
% \hfill
% \DatestampYMD{2022}{11}{19}
% \begin{Detail}
% \Item
% Awarded at the Tenth Annual Chess Tournament held during Open House.
% \end{Detail}

\Gap
\BulletItem
\textbf{The GBA International Programming Contest: }Third Prize
\hfill
\DatestampYMD{2025}{4}{19}
\begin{Detail}
\Item
The GBA international programming contest is organized by the Shenzhen SLAI and The CUHKSZ, provides a
platform for students to demonstrate their
innovative abilities and skills in analyzing and solving problems.
\end{Detail}

\Gap
\BulletItem
\href{http://www.cmathc.cn/}{\textbf{The Chinese Mathematics Competitions}}: Provincial First Prize
\hfill
\DatestampYM{2021}{12}
% \newline
% Global Science, Technology, Engineering, and Mathematics Foundation
\begin{Detail}
\Item
National calculus competition for undergraduates. Obtained in the first semester of university, final standing 317/2300
\end{Detail}

\Gap
\BulletItem
\href{https://www.mcm.edu.cn/html_cn/node/5b6ed02462b2ea3d79206e445af3433a.html}{\textbf{Contemporary Undergraduate Mathematical Contest in Modeling}}: National Second Prize
\hfill
\DatestampYM{2022}{10}
% \newline
% Global Science, Technology, Engineering, and Mathematics Foundation
\begin{Detail}
\Item
One of 1,445 winning teams nationwide
\end{Detail}

\Gap
\BulletItem
\href{www.comap.com}{Interdisciplinary Contest in Modeling}: Meritorious Winner
\hfill
\DatestampYM{2024}{02}
% \newline
% Global Science, Technology, Engineering, and Mathematics Foundation
% \begin{Detail}
% \Item
% One of 1,445 winning teams nationwide
% \end{Detail}

\Gap
\BulletItem
Collegiate Computer Systems \& Programming Contest: Bronze Medal
\hfill
\DatestampYM{2022}{12}
% \newline
% Global Science, Technology, Engineering, and Mathematics Foundation
% \begin{Detail}
% \Item
% One of 1,445 winning teams nationwide
% \end{Detail}

\Gap
\BulletItem
First-Class Scholarship for Outstanding Students of the School
\hfill
\DatestampYM{2022}{5}
% \newline
% Global Science, Technology, Engineering, and Mathematics Foundation
% \begin{Detail}
% \Item
% One of 1,445 winning teams nationwide
% \end{Detail}

\Gap
\BulletItem
Zhou Peiyuan Mechanics Competition for College Students: Provincial Second Prize
\hfill

\Gap
\BulletItem
National Undergraduate Electronics Design Contest: Provincial Third Prize
\hfill
%%%%%%%%%%%%%%%%%%%%%%%%%%%%%%%%%%%%%%%%%%%%
%% Organization Experience %%
%%%%%%%%%%%%%%%%%%%%%%%%%%%%%%%%%%%%%%%%%%%%

\Section
{Research\newline
Experience}
{Professional Affiliations \& Activities}
{PDF:ProfessionalAffiliationsActivities}

\Entry
\textbf{Harbin Engineering University ACM/ICPC Team}
\hfill
\DatestampYM{2022}{03} --
\DatestampYM{2022}{12}
\newline
Represent the school in ACM/ICPC competitions, be responsible for relevant training courses within the school, and organize regional competitions.

\BigGap
\textbf{Harbin Engineering University Aircraft Control Laboratory}
\hfill
\DatestampYM{2022}{12} --
\DatestampYM{2023}{11}
\newline
Mainly research aircraft simulation, theory and application of advanced control algorithms, and application of statistical learning methods in the control field.
\begin{Detail}
\BulletItem
Built a high-performance C++/Python-mixed aircraft simulation for the lab.
\BulletItem
Solve problems in discrete mathematics for the lab.
%\BulletItem
%Represented the college to participate in a research institute's cooperation projects and was responsible for the implementation of reinforcement learning algorithms.
\end{Detail}

\BigGap
\textbf{Academy of Aerospace Solid Propulsion Technology}
\hfill
\DatestampYM{2023}{7} --
\DatestampYM{2023}{10}
\newline
Build a reinforcement learning environment framework and write high-performance simulation of aircraft aerodynamics. Training multi-agent reinforcement learning algorithms (MADDPG) in engineering projects.

\BigGap
\href{http://rl.beiyang.ren/}{\textbf{Tianjin University Deep Reinforcement Learning Laboratory}}
\hfill
\DatestampYM{2023}{11} -- Present
\newline
A top research team in the field of reinforcement learning in China, focusing on cutting-edge research in the field of decision-making intelligence.
\begin{Detail}
\BulletItem
Completed a paper as co-author and submitted to \textit{AAAI}.
\BulletItem
As a core developer, formed a team to participate in a NeurIPS competition and won a top prize.
\end{Detail}

%%%%%%%%%%%%%%%%%%%%%%%%%
%% Engineering Projects %%
%%%%%%%%%%%%%%%%%%%%%%%%%

\clearpage
\Section
{Engineering Projects}
{Engineering Projects}
{PDF:ResearchExperience}

\Entry
\textbf{Swarm Aircraft Decision Intelligence Project}
\hfill
\DatestampYM{2023}{07}
\newline
Academy of Aerospace Solid Propulsion Technology
\BulletItem
Programming aircraft six degrees of freedom pneumatic physics simulation
\BulletItem
Responsible for the design, deployment and training of reinforcement learning algorithms
\BulletItem
Coordinate with other departments to complete joint programming

\BigGap
\Entry
\textbf{Aerodrome: A Lightweight physics simulation platform for reinforcement learning}
\hfill
\DatestampYM{2025}{01}
\newline
C++/Python joint programming reinforcement learning platform, my undergraduate graduation project
\BulletItem
High-performance C++ aircraft dynamics simulation
\SubBulletItem
Highly extensible code framework
\SubBulletItem
Standardized reinforcement learning environment and Python interface

Github Repo: \url{https://github.com/CH4ACKO3/Aerodrome}

\BigGap
\Entry
\textbf{ResearchFlow: ML research workflow management toolkit}
\hfill
\DatestampYM{2025}{10}
\newline
Python-based research toolkit for flexible process scheduling, model/data storage and management
\Gap
\BulletItem
JSON-based indexed file storage
\SubBulletItem
Retrieve model and result files by metadata querying, with several training parameter
\SubBulletItem
Multi-process and crash-safe
\BulletItem
Asynchronous multi-GPU task scheduler
\SubBulletItem
Dynamic computing resource monitoring and online GPU scheduling
\SubSubBulletItem
Automatically suspend tasks when the system is critical to avoid crashes.
\SubSubBulletItem
Withdraw tasks from a GPU and reassign at any time.
\SubSubBulletItem
Fully asynchronous and low-latency file I/O, task pool management, compute device management
\SubBulletItem
Interact with original text documents, researcher-friendly
\SubBulletItem
Robust crash and interrupt handling
\SubBulletItem
Comprehensive test suite with 54 test cases (unit, integration, concurrency)

Github Repo: \url{https://github.com/CH4ACKO3/researchflow}
%%%%%%%%%%%%%%%%%%
%% PUBLICATIONS %%
%%%%%%%%%%%%%%%%%%

\Section
{Publications}
{Publications}
{PDF:Publications}

\Entry
\href{https://arxiv.org/abs/2408.15501}{MODULI: Unlocking Preference Generalization via Diffusion Models for Offline Multi-Objective Reinforcement Learning}. (\textit{ICML2025})

Yifu Yuan, \textbf{Zhenrui Zheng}, Zibin Dong, Jianye HAO.

Github Repo: \url{https://github.com/pickxiguapi/MODULI}

\end{Body}

%%%%%%%%%%%
% CV NOTE %
%%%%%%%%%%%

\BigGap
\UseNoteFont%
\null\hfill%
[\textit{\CVNote}]

\end{document}
